% YY13100B
% Standardmaßnahmen bei vital bedrohten Patienten
% 14.0
% 2013-11-30
\label{m:standardmassnahmen-vital}%
\label{topic:standardmassnahmen-vital}%
\label{proc:YY13100B}%
\begin{enumerate}
  \item \EE{Situationsgerechte Lagerung}
  \item \EE{Sauerstoffgabe}
  (\prettyref{m:sauerstoffberieselung}): je nach Indikation,
  allgemeiner Zielwert: Sp\ce{O2} von \VonBis{94}{98}\PC*
  \item Ggfs. \EE{Notarzt-Nachforderung} :
  mit kurzer Begründung
  
  \Ca{Bei manchen Notfällen keine Notarztnachforderung wenn
  der Patient zeitkritisch ist, siehe jeweilige spezielle
  Maßnahmen!}
  \item \EE{Engmaschiges, bestmögliches Monitoring}: Je nach
  vorhandenem Material
  
  RR, HF, Pulsoxymetrie, EKG,
  Sitzwache/Patientenbeobachtung, etc.\footnote{Zuerst werden die Patientenbeobachtung
  und -- wenn vorhanden -- die
  Pulsoxymetrie eingesetzt. Überwachungsgeräte, deren Anlage
  zeitintensiv ist (EKG, \ldots), sollen erst dann verwendet
  werden, wenn alle dringlicheren Maßnahmen durchgeführt und
  der Einschätzungsblock beendet ist! Die
  Patientenbeobachtung bleibt immer ein wesentlicher Teil
  des Monitorings!}
  \item \EE{Reanimationsbereitschaft} herstellen
  (\prettyref{m:reanimationsbereitschaft})
\end{enumerate}

{\scriptsize(Reihenfolge zählt!)}